% !PRINTED 524
\weberpage{524}

%<<CONT
wo für $\gamma_2(\omega)$ und $j(\omega)$ die Werte (5) einzusetzen sind. Diese Gleichungen lassen sich noch zerlegen, und man erhält für $x$ selbst viel einfachere Gleichungen:
\weq{8}{\begin{array}{ll}
f(\sqrt{-51})^3 = 2x, & x^3 - 4x^2 - x - 1 = \sqrt{17}\,x^2,\\
f(\sqrt{-91}) = x, & x^3 - 2x^2 - x - 2 = \sqrt{13}\,x,\\
f(\sqrt{-99})^3 = 2x, & x^3 - 13x^2 - 4x - 1 = \sqrt{33}\left(2x^2 + x\right).
\end{array}}
Die Darstellung der Klasseninvarianten $j(\omega)$ durch eine einzige Quadratwurzel ist immer dann möglich, wenn zwei Geschlechter und in jedem Geschlecht eine Klasse vorhanden ist. Diese Werte von $m$ finden sich in der \gesperrt{Gauss}schen Tafel der Klassenzahlen (Werke, Bd.\,II, S.\,450\,ff.) unter der Bezeichnung II, 3, von denen die auszuwählenden sind, die $\equiv 3\ (\mathrm{mod}\ 8)$ sind. Ihre Anzahl ist aller Wahrscheinlichkeit nach endlich, wie die erwähnte Tafel aufweist; es sind die Zahlen:
\[m = 35, \quad 51, \quad 75, \quad 91, \quad 99, \quad 115, \quad 123,\]
\[147, \quad 187, \quad 235, \quad 267, \quad 403, \quad 427.\]
Wenn die Formenklassen in eine beliebige Anzahl von Geschlechtern zerfallen, aber in jedem Geschlecht nur eine Klasse enthalten ist, dann lassen sich nach dem in § 138 bewiesenen Satze alle Klasseninvarianten durch \gesperrt{Quadratwurzeln} ausdrücken. Von \gesperrt{Euler} und \gesperrt{Gauss} ist durch Induktion geschlossen, daß es nur eine endliche Anzahl, nämlich 65, solcher Diskriminanten gibt\footnote{\gesperrt{Euler}, Nouv. Mém. de Berlin 1776, S.\,338. \gesperrt{Gauss}, Disq. art. 303. \gesperrt{Euler} ist auf diese Zahlen auf einem anderen Wege gelangt, nämlich von der Aufgabe, große Primzahlen zu ermitteln (Vgl. die Straßburger Dissertation von \gesperrt{Peter Meyer}: Beweis eines von \gesperrt{Euler} entdeckten Satzes, betreffend die Bestimmung von Primzahlen, Straßburg 1906). Diese Zahlen sind:

1, 2, 3, 4, 5, 6, 7, 8, 9, 10, 12, 13, 15, 16, 18, 21, 22, 24, 25, 28, 30, 33, 37, 40, 42, 45, 48, 57, 58, 60, 70, 72, 78, 85, 88, 93, 102, 105, 112, 120, 130, 133, 165, 168, 177, 190, 210, 232, 240, 253, 273, 280, 312, 330, 345, 357, 385, 408, 462, 520, 760, 840, 1320, 1365, 1848.

% Source chapter-closing rule follows the completed footnote.
\par\medskip\noindent\makebox[\linewidth]{\rule{0.25\linewidth}{0.2pt}}}. In der Abhandlung \glqq Zur komplexen Multiplikation elliptischer Funktionen\grqq{} (Mathematische Annalen, Bd.\,33) habe ich diese Zahlen alle berechnet. Sie sind auch in der Tafel der Klasseninvarianten, die diesem Buche beigefügt ist, enthalten.
